% prefacio.tex
%
% Copyright (C) 2022--2026 José A. Navarro Ramón <janr.devel@gmail.com>
% Licencia del código GPLv2
% Licencia Creative Commons Recognition Non-Commercial Share-alike.
% (CC-BY-NC-SA)

\chapter{Prefacio}
Objetivos que pretendo conseguir:
\begin{enumerate}
\item Reunir en unos apuntes diversos conceptos, desarrollos y ejercicios matemáticos
  para disponer de un texto de consulta sobre esta disciplina, fundamentalmente aunque
  no necesariamente, como apoyo al estudio de la física.
\item Explicar a \emph{mi yo del futuro} los conceptos y técnicas anteriores, para que
  pueda volver a asimilarlas en la medida de lo posible, si lo llegara a necesitar.
\end{enumerate}

No pretendo formar un compendio completo, ni siquiera ordenado, sobre las matemáticas
aplicadas a la física.
Por tanto, he añadido aquí los detalles matemáticos que he necesitado o echado en falta
en algún momento de mi estudio, principalmente de física, pero no limitado a esta
materia.

La motivación de estos apuntes ha sido para explicármelo a mí.
Se puede desprender de todo lo anterior que he hecho hincapié en todo aquello que en
algún momento he encontrado más complicado, y no he dedicado ninguna o casi ninguna
atención a detalles que tengo suficientemente asimilados.

Me pregunto si estos apuntes podrían servir a alguien. No lo sé. Para ello, se debería
tener un nivel similar al mio, pues si fuera superior podría considerarse muy elemental
y una pérdida de tiempo. Si fuera inferior, se necesitaría profundizar en otros conceptos
que doy por sabidos.

En todo caso, según la licencia, cualquiera que pudiera estar interesado en el texto y no
lo considerara lo suficientemente didáctico o accesible, puede modificarlo y usarlo para
sus propios objetivos. Para ello, consulta la licencia, que he separado en dos partes:
\begin{enumerate}
\item El código fuente, con licencia
  \href{https://www.gnu.org/licenses/old-licenses/gpl-2.0.html#SEC1}{GPLv2}, se encuentra
  en \href{https://github.com/janr57/retazosmatematicas}{GitHub}.
\item En cambio, el fichero en formato de texto obtenido, que se puede distribuir por
  separado, tiene licencia
  \href{https://creativecommons.org/licenses/by-nc-sa/4.0/}{Creative Commons BY-NC-SA}.
\end{enumerate}

 
%%% Local Variables:
%%% mode: latex
%%% TeX-engine: luatex
%%% TeX-master: "../retazosmatematicas.tex"
%%% End:
